\documentclass[11pt]{article}
\usepackage{rrs}

\title{Introduction to Root Rewrite Systems}
\author{Michal Szabados}
\date{\today}

\begin{document}
\maketitle

Root Rewrite System (RRS) is a computation model operating on \emph{terms}. The
definitions below use ideas from term rewriting systems and functional
language features in e.g. ML, Haskell or Rust.

\section{Signatures and Terms}

Definitions in this section follow the term rewriting systems theory.

Throughout, we fix a countably infinite set $\Variables$ of \emph{variables},
disjoint from any other set we consider (in particular from the constructors of
any signature). As a convention, $x, y, z \in \Variables$.

\begin{definition}[Signature]
A \emph{signature} is a tuple $(\Constructors, \arity)$ where:
\begin{itemize}
  \item $\Constructors$ is a finite set;
  \item $\arity \colon \Constructors \to \N$ is a function.
\end{itemize}
Constructors of arity zero are called \emph{constants}, of arity one
\emph{unary} and of arity two \emph{binary}.
\end{definition}

\begin{definition}[Terms]
Let $(\Constructors, \arity)$ be a given signature. For $V \subseteq \Variables$ define $\Terms(V)$, the set of
\emph{terms with variables in $V$,} as the smallest set such that:
\begin{itemize}
  \item $V \subseteq \Terms(V)$;
  \item For a constant $C \in \Constructors$ we have $C \in \Terms(V)$;
  \item For $C \in \Constructors$ with $n := \arity(C) > 0$ and $t_1, \dots, t_n \in \Terms(V)$
        we have $C(t_1, \dots, t_n) \in \Terms(V)$.
\end{itemize}
\end{definition}

We understand terms as trees rather than strings. In particular, rooted trees with ordered
children. The leaves are variables and constants, and inner nodes are constructors
with as many children as is their arity.

\begin{definition}[Term variables]
Let a signature be given and $t \in \Terms(\Variables)$. Define $\Var(t)$, the
\emph{variables occurring in $t$}, as the smallest $V \subseteq \Variables$ such that
$t \in \Terms(V)$. If $V$ is the empty set we call $t$ \emph{ground}.

The set of
all ground terms is denoted simply $\Terms$, and the set of all terms with
variables, $\Terms(\Variables)$, is denoted $\TermsVar$.
\end{definition}

\begin{example}[Natural numbers]\label{ex:natural-numbers}
Let us construct terms representing natural numbers. We define the signature by:
\begin{align*}
  \Constructors &= \{\rrsname{Zero}, \rrsname{Succ}\} \\
  \arity \colon \Constructors &\to \N \\
  \rrsname{Zero} &\mapsto 0 \\
  \rrsname{Succ} &\mapsto 1
\end{align*}
Example terms representing 0, 1, 2 are:
\begin{itemize}
  \item $\rrsname{Zero}$
  \item $\rrsname{Succ}(\rrsname{Zero})$
  \item $\rrsname{Succ}(\rrsname{Succ}(\rrsname{Zero}))$
\end{itemize}
These are ground, and all are elements of $\Terms$. An example element of $\TermsVar$
is $\rrsname{Succ}(\rrsname{Succ}(x))$.
\end{example}

\begin{definition}[Substitution]
Let a signature be given. A map $\sigma \colon \Variables \to \Terms$ is
called \emph{a substitution.} We naturally extend $\sigma$ to a map
$\hat{\sigma} \colon \TermsVar \to \Terms$ as follows:
\begin{itemize}
  \item For $x \in \Variables$ we have $\hat{\sigma}(x) := \sigma(x)$;
  \item For a constant $C \in \Constructors$ we have $\hat{\sigma}(C) := C$;
  \item For $C \in \Constructors$ with $n := \arity(C) > 0$ and $t_1, \dots, t_n \in \TermsVar$
        we have $\hat{\sigma}(C(t_1, \dots, t_n)) := C(\hat{\sigma}(t_1), \dots, \hat{\sigma}(t_n))$.
\end{itemize}
\end{definition}

\section{Root Rewrite Systems}

Now that we have the definition of terms and related concepts, the definition
of a root rewrite system is rather straightforward.

\begin{definition}[Rewrite rule]
Let a signature be given. A \emph{rewrite rule} is a tuple
$(\tleft, \tright) \in \TermsVar^2$ such that:
\begin{itemize}
  \item every variable in $\tleft$ occurs at most once;
  \item every variable in $\tright$ occurs at most once;
  \item every variable in $\tright$ occurs in $\tleft$.
\end{itemize}
A rewrite rule \emph{matches} a term $t \in \Terms$ if there exists a substitution $\sigma$
such that $\hat{\sigma}(\tleft) = t$. In that case, $\hat{\sigma}(\tright)$ does not depend
on the choice of $\sigma$ and we call it \emph{$t$ rewritten by the rule}.
\end{definition}

Note that the rules do not allow for duplication of variables. Subterms bound by variables can
only be moved around or discarded.

\begin{definition}[Root Rewrite System]
A \emph{Root Rewrite System} is a tuple $\rrsname{RRS} = (\Sig, \rrsname{Rules})$
where $\Sig$ is a signature and $\rrsname{Rules}$ is an ordered list of rules.

Given an input term $t_0 \in \Terms$, in one \emph{step} the machine rewrites
$t_i \to t_{i+1}$ using the first matching rule. The sequence
terminates at $t_n$ if no rule matches; then $t_n$ is the output of the machine.
If the sequence does not terminate the machine does not produce an output.
\end{definition}

\begin{note}
In term rewriting terminology this is a \emph{priority root-rewrite system}: rules apply
only at the root, and the first matching rule wins. Moreover, the rules are \emph{linear.}
Note that $t_n$ defined above is the \emph{normal form} of $t_0$.
\end{note}

We define a simple programming language to give examples of RSSes. A program is of the form:

\begin{rrscode}
rule <lhs1> -> <rhs1>
rule <lhs2> -> <rhs2>
...
\end{rrscode}

Every line defines a rewrite rule, ordered from top to bottom. The signature is implicitly
defined as the set of all constructors occurring in the rules, with arity given by their
usage (which must be consistent). To distinguish constructors from variables, we require that
constructors start with an uppercase letter and variables with a lowercase one.

\begin{example}[Addition of natural numbers]
With the natural number definition from Example \ref{ex:natural-numbers}, let us add an
$\rrsname{Add}$ constructor with arity two:

\begin{rrscode}
rule Add(Succ(x), y) -> Add(x, Succ(y))
rule Add(Zero, y) -> y
\end{rrscode}

The idea is to move all $\rrsname{Succ}$ constructors from the first child to the second
one-by-one. I.e., $\rrsname{Add}(x+1, y)$ rewrites to $\rrsname{Add}(x, y+1)$. A computation
of 2+1 would look like this:
\begin{align*}
  &\rrsname{Add}(\rrsname{Succ}(\rrsname{Succ}(\rrsname{Zero})), \rrsname{Succ}(\rrsname{Zero})) \\
  {}\to{} &\rrsname{Add}(\rrsname{Succ}(\rrsname{Zero}), \rrsname{Succ}(\rrsname{Succ}(\rrsname{Zero}))) \\
  {}\to{} &\rrsname{Add}(\rrsname{Zero}, \rrsname{Succ}(\rrsname{Succ}(\rrsname{Succ}(\rrsname{Zero})))) \\
  {}\to{} &\rrsname{Succ}(\rrsname{Succ}(\rrsname{Succ}(\rrsname{Zero})))
\end{align*}
which is the term representing 3, as expected.

\end{example}

\begin{example}[Reverse a list]
TODO
\end{example}


\end{document}
